\documentclass[a4paper,12pt]{article}

% --- Pakete ---
\usepackage[ngerman]{babel}
\usepackage[utf8x]{inputenc}
\usepackage[T1]{fontenc}
\usepackage{eurosym}
\usepackage{amsmath,amssymb}
\usepackage{listings}
\usepackage{xcolor}

% --- Python Code Stil ---
\lstset{
	language=Python,
	basicstyle=\ttfamily\small,
	backgroundcolor=\color{gray!10},
	frame=single,
	numbers=left,
	numberstyle=\tiny,
	breaklines=true,
	showstringspaces=false
}

\title{Übungsblatt 1}
\author{}
\date{}

\begin{document}
	\maketitle
	
	\section*{Aufgabe 1.1}
	\subsection*{Diskret/Stetig (Quantitativ), Ordinal/Nominal (Qualitativ)}
	
	\begin{enumerate}
		\item Anzahl der Insassen in einem PKW bei der Verkehrszählung\\
		\textbf{Diskret} (abzählbar, nur ganze Zahlen)
		
		\item Höchstgeschwindigkeit eines PKWs\\
		\textbf{Stetig} (überabzählbar, kann auch Gleitkommazahlen beinhalten)
		
		\item Schultypen\\
		\textbf{Nominal} (keine Zahl)
		
		\item Temperaturangaben in °C\\
		\textbf{Stetig} (Kommazahlen)
		
		\item Fassungsvermögen einer Festplatte in MB\\
		\textbf{Diskret} (halbe Byte sind nicht möglich; heutige Hardware kann keine unendlichen Nachkommastellen darstellen)
		
		\item Zugriffszeiten auf die Daten einer Festplatte in Sekunden\\
		\textbf{Diskret} (keine halben Sekunden)
		
		\item Energieeffizienz von Elektrogeräten\\
		\textbf{Ordinal} (nach Rangordnung sortierbar)
		
		\item Trinkmenge eines Menschen pro Tag in Litern\\
		\textbf{Stetig} (in einem Intervall $(0, \text{max. mögliche Menge})$)
		
		\item Anzahl Likes eines Instagram-Posts\\
		\textbf{Diskret} (nur ganze Likes)
		
		\item Nationalität eines Menschen\\
		\textbf{Nominal} (nicht ordnenbar)
		
		\item Anfangsbuchstabe des Vornamens eines Menschen\\
		\textbf{Kommt darauf an:}\\
		Wenn als 1-Byte-Zahl (UTF-8/ASCII) dargestellt → \textbf{Ordinal}.\\
		Wenn als Buchstabe betrachtet → \textbf{Nominal}.
		
		\item Wohnort eines Menschen\\
		\textbf{Nominal} (String)
	\end{enumerate}
	
	\section*{Aufgabe 1.2}
	\subsection*{Aufgabe 1.2.1}
	
	\begin{enumerate}
		\item \textbf{Statistische Einheiten:} Anzahl an Aufrufen
		\item \textbf{Grundgesamtheit:} Demographie und Nationalität der Zuschauer
		\item \textbf{Stichprobe:} Kinder im Alter von 6–10 Jahren
		\item \textbf{Merkmale:}
		\begin{enumerate}
			\item CPM eines Videos
			\item Uhrzeiten, bei denen das Video geschaut wird
			\item Videolänge
		\end{enumerate}
		\item \textbf{Beispielhafte Merkmalsausprägungen:}
		\begin{enumerate}
			\item 3{,}45 €
			\item 18:00 Uhr
			\item 12 Minuten
		\end{enumerate}
		\item \textbf{Eigenschaften der Merkmale:}
		\begin{enumerate}
			\item Quantitativ, stetig
			\item Quantitativ, diskret
			\item Quantitativ, diskret
		\end{enumerate}
	\end{enumerate}
	
	\subsection*{Aufgabe 1.2.2}
	
	Die durchschnittliche Watchtime von Zuschauern basierend auf den im Video getätigten Aktionen, 
	da sie viel darüber aussagt, wie interessant der Content ist und ob man ihn anpassen sollte.
	
	\section*{Aufgabe 1.3}
	
	\subsection*{a)}
	
	\begin{lstlisting}
		import numpy as np
		
		cash = {5, 10, 20, 50, 100, 200, 500}
		print("Median:", np.median(list(cash)))  # 50
	\end{lstlisting}
	
	\subsection*{b)}
	
	\begin{lstlisting}
		cash.add(1000)
		print("Median mit hinzugefuegtem 1000 Euro-Schein:", np.median(cash)) # 75
	\end{lstlisting}
	
\end{document}
